\answer{ex:08:1}{(а)~$\approx1.7\cdot10^6$. (б)~$g=\frac{0.9}{\sqrt{0.19}}\,\sigma_{\text{после}}/\sigma_{\text{до}}\approx16.5$: ответ зависит только от отношения шумов.}
\answer{ex:08:2}{(а)~Собственные числа $-(1\pm g\beta)/\tau$; разность затухает за $\tau/(1-g\beta)$: $40\ms$ при $g\beta=0.5$ (разница входов усилена втрое), $200\ms$ при $0.9$ (в $19$ раз). (б)~$0.905$ и $1.105$; между границами побеждает только сильный вход.}
\answer{ex:08:3}{$P(k)=e^{gu_k/\sigma}/\sum_le^{gu_l/\sigma}$, то есть~\eqref{eq:softmax}; $g=10\ln9\approx22$.}
\answer{ex:08:4}{$\bar\Delta=1/3$, $g^*\approx3\ln(1/h)$: $21$, $14$, $9$ против $16$--$23$, $11$ и $8$ в модели.}
\answer{ex:08:5}{$g^*$ от $16$--$23$ при $h=0.001$ до $6$--$8$ при $h=0.2$; оценка $3\ln(1/h)$ верна в пределах полутора раз при $h\le0.05$; при $h\gtrsim\eta/10$ правило $Q\leftarrow Q+\eta(R-Q)$ не успевает усвоить добытое поиском, и $g^*$ застревает около $8$.}
\answer{ex:08:6}{$T_p=\tau\ln9=44\ms$ при $g_{\text{lo}}=0$ (модель $44\ms$) и $70\ms$ при $g_{\text{lo}}=0.25$: пачка в $2$--$3\tau$ при усилении около нуля.}
\answer{ex:08:7}{(а)~Лучшее постоянное $0.925$ ($g=8$), оракул $0.929$ ($16/8$) --- выигрывать почти нечего. (б)~Например, спокойные эпохи с редкими наградами ($p\in[0,0.3]$) и переменчивые с $h=0.1$: $0.850$ против $0.872$; подстройка забирает около четверти при $\eta=0.2$ и почти всё при $\eta=0.05$.}
